problem:  
Let \((X,L)\) be a polarized compact Kähler manifold with nontrivial connected automorphism group \(G = \mathrm{Aut}^0(X,L)\), and let \(K \subset G\) be a maximal compact subgroup. Assume that \(X\) admits an extremal Kähler metric \(\omega_{\mathrm{ext}}\) with extremal vector field \(\xi \in \mathrm{Lie}(K)\). For each integer \(k>0\), consider the normalized Bergman embedding \(\Phi_k: X \to \mathbb{P}(H^0(X,L^k)^*)\) determined by a \(K\)-invariant orthonormal basis of sections for \(\omega_{\mathrm{ext}}\), and let \(\omega_k = \frac{1}{k}\Phi_k^*(\omega_{\mathrm{FS}})\) be the induced metric in \(c_1(L)\).  

Define the **relative center of mass** \(\mu_k \in \mathrm{Lie}(K)\) of \(\Phi_k\) as the component of the Bergman moment map orthogonal to the infinitesimal automorphisms generated by \(\xi\).  

Is it true that \(\|\mu_k\| = O(k^{-\alpha})\) for some \(\alpha>0\) uniformly in \(k\)? Equivalently, does the sequence \((\omega_k)\) asymptotically eliminate all Hamiltonian drift orthogonal to the extremal vector field at a polynomial rate?  
If so, can the rate of decay of \(\|\mu_k\|\) serve as an *analytic invariant* distinguishing **uniform relative K-stability** from mere relative K-polystability?  

Why is it a "good" problem:  
This formulation isolates a sharp and currently unresolved question at the analytic–algebraic interface of extremal Kähler geometry. While existence and weak convergence of relative balanced metrics have been studied, quantitative asymptotics of the moment map’s relative component remain uncharted. Determining whether the center of mass decays polynomially — and how its rate encodes stability strength — would introduce a new bridge between Geometric Invariant Theory notions (uniform vs. polystability) and the asymptotic behavior of Bergman-type metrics. The problem is precise, well-posed, and simple to state, yet it touches a profound open frontier: developing effective analytic measures of stability in the equivariant (relative) setting beyond existence statements. This line of inquiry could yield new quantitative refinements of the Yau–Tian–Donaldson philosophy for manifolds with symmetries, making it a "good" and forward-looking research problem.